\section{The invariant and the main theorem}
\label{sec:main}
For a full-dimensional convex body $K\subset\R^d$, let $|K|$ denote its
Euclidean volume. For a \emph{unit} vector $u$, let $\pi_K(u)$ be the
$(d-1)$-dimensional volume of its orthogonal projection onto $u^\perp$.
We reserve the homogeneous notation
\[
 b_K(z)=\begin{cases}\norm z\pi_K(z/\norm z),&z\ne0,\\0,&z=0\end{cases}
\]
for nonunit vectors. The projection body $\Pi K$ has support function
$h_{\Pi K}=b_K$. If $L$ has dimension $p$, set
\[
 R(L)=\frac{|\Pi L|}{|L|^{p-1}},\qquad
 \pyr L=\conv\bigl((L,0)\cup\{(0,1)\}\bigr).
\]
We use exactly the projection-body pyramid invariant of the retained
Entry 005 project:
\begin{equation}\label{eq:actualdefect}
 a(K)=\left(\frac d{d+1}\right)^d
       \frac{R(\pyr K)}{R(K)}-1,
 \qquad e(K)=a(K)-\frac1{d+1}.
\end{equation}
For a full-dimensional simplex $S\subset K$ with centroid $g_S$, define
\begin{equation}\label{eq:excess}
 E(K,S)=\inf\{\tau\ge0:K\subset g_S+(1+\tau)(S-g_S)\}.
\end{equation}
The dilation factor itself is $1+E$. The simplex in this definition is
prescribed; neither the simplex nor its center is optimized in $E(K,S)$.

\begin{theorem}[Sharp vertex-excess exponent]\label{thm:main}
Let $d\ge3$ and $r\ge1$ be integers, and put
\[
 \begin{gathered}
 n=d+1,\qquad k=\min(r,d-1),\qquad H_d=3d^2(d+1)^3(d+2),\\
 T_{d,r}=\min\{64(d+1)r,16(d+1)^2\}.
 \end{gathered}
\]
For every full-dimensional convex $d$-polytope $K$ with at most $d+1+r$
vertices, and every prescribed maximum-volume inscribed simplex $S$,
\begin{equation}\label{eq:main}
 E(K,S)\le T_{d,r}H_d^{1/k}e(K)^{1/k}.
\end{equation}
The invariant $e(K)$ is nonnegative. The exponent $1/k$ is optimal:
there is an explicit family $(K_t,S_t)$ in this same class, with $S_t$
maximum and $e(K_t)\to0$ as $t\downarrow0$, such that
\begin{equation}\label{eq:sharp_limit}
 \lim_{t\downarrow0}\frac{E(K_t,S_t)}{e(K_t)^{1/k}}
 =n\left(\frac{n^2}{k(k+1)}\right)^{1/k}.
\end{equation}
Consequently no exponent $\beta>1/k$ admits a uniform bound near zero
defect, even with an arbitrary finite coefficient depending on $d,r$.
\end{theorem}

One extra vertex gives a linear defect estimate. Two extra vertices give
exponent $1/2$ in every dimension $d\ge3$. Saturation at the unrestricted
exponent $1/(d-1)$ occurs when the vertex allowance reaches $2d$.
The coefficient in \eqref{eq:main} is explicit but is not asserted to be
optimal. The case $r=0$ is trivial: a full-dimensional polytope with at most
$d+1$ vertices is a simplex. The theorem concerns $d\ge3$; the two-dimensional
body used in the sharpness construction is computed explicitly below.

The proof separates three geometric operations. First, a finite apex count
controls the volume missing from a cone and survives a ray-collapsing
projection. Second, actual facet normals yield an enclosing simplex with a
relative projection error of at most $H_de(K)$. Third, a near-vertex
condition for that enclosing simplex controls the originally prescribed
maximum, including maxima whose vertices are not vertices of $K$.

\subsection{Context and scope}
Mixed-volume polynomials, their monotonicity, the first mixed-volume formula,
and Minkowski's first inequality are classical inputs
\cite{SchneiderBook,BS}. They are not new ingredients of this work.
Simplex stability for other functionals is studied, for example, by
Schneider~\cite{SchneiderStability}; cone-volume stability for subspace
concentration and the $U$-functional is treated in~\cite{BH}. Projection-body
volume questions also have substantial earlier literature
\cite{Brannen,Saroglou}. These neighboring statements involve different
functionals or classes and are not used to infer a novelty claim here.
The source comparison accompanying this manuscript is focused rather than
exhaustive; in particular it does not certify the complete overlap with
all earlier projection-body results.

The cone count, projection collapse, and actual-facet specialization are
the added steps relative to the retained unrestricted proof. For transparency,
we also reproduce finite versions of its determinant assignment and
retention arguments. Thus the mathematical chain below can be read without
running the earlier formal source tree.
